\SourcePage{278}
\SGAsection{7}{对固有平坦群胚取商}
\begin{SGAresult}{\Term{theorem}{定理}}{7.1}
\footnote{编者注：这里指出 S. Keel 与 S. Mori 的文章 [KM97]，其中证明了下述\Term{theorem}{定理}。设 $X$ 是在\Term{locally-noetherian}{局部诺特}底\Term{scheme}{概形} $S$ 上的\Term{finite-type}{有限型}\Term{algebraic-space}{代数空间}，$j:R\to X\times_SX$ 是\Term{flat}{平坦}\Term{groupoid}{群胚}，其\Term{stabilizer}{稳定子} $j^{-1}(\Delta_X)$ 在 $X$ 上\Term{finite-morphism}{有限}；则存在一个\Term{algebraic-space}{代数空间}，它既是 $X$ 对 $R$ 的\Term{geometric-quotient}{几何商}，又是\Term{uniform-categorical-quotient}{一致范畴商}（quotient catégorique uniforme）；而且，若 $j$ 是\Term{separated}{分离}的，这个\Term{quotient}{商}也\Term{separated}{分离}。特别地，若\Term{flat}{平坦} $S$-\Term{group-scheme}{群概形} $G$ 在 $X$ 上\Term{proper-action}{固有地作用}，且\Term{stabilizer}{稳定子}\Term{finite-morphism}{有限}（即\Term{morphism}{态射} $G\times_SX\to X\times_SX$，$(g,x)\mapsto(x,g\cdot x)$，是\Term{proper}{固有}的，且\Term{diagonal}{对角线}的\Term{stabilizer}{稳定子}在 $X$ 上\Term{finite-morphism}{有限}），则存在\Term{geometric-quotient}{几何商} $X\to X/G$。当 $G$ 是\Term{reductive-group-scheme}{约化群概形}时，这是 J. Kollár 的一个结果 [Ko97]。}
设 $S$ 是\Term{locally-noetherian}{局部诺特}\Term{scheme}{概形}，且
\[
\begin{tikzcd}
X_2\arrow[r,shift left=1.1ex,"d'_2"]\arrow[r,"d'_1" description]\arrow[r,shift right=1.1ex,"d'_0"']&X_1\arrow[r,shift left=.6ex,"d_1"]\arrow[r,shift right=.6ex,"d_0"']&X_0
\end{tikzcd}
\]
是一个 $(\mathrm{Sch}/S)$-\Term{groupoid}{群胚}，其中 $d_1$ 是\Term{proper}{固有}\Term{flat}{平坦}的，$X_0$ 在 $S$ 上\Term{quasi-projective}{拟射影}，\footnote{编者注：对 $X_0$ 的这个假设是必要的，参见\Term{theorem}{定理} 4.1 中的编者注 (17)。}且\Term{set-component}{分量}为 $d_0,d_1$ 的\Term{morphism}{态射} $d:X_1\to X_0\times_SX_0$ 是\Term{quasi-finite}{拟有限}的。那么：

(i) $(d_0,d_1)$ 在 $(\mathrm{Sch}/S)$ 中存在\Term{coequalizer}{余核} $(Y,p)$；而且，这样的 $(Y,p)$ 也是 $(d_0,d_1)$ 在全体\Term{ringed-space}{赋环空间}的\Term{category}{范畴}中的\Term{coequalizer}{余核}。

(ii) $p$ 为\Term{surjectivity}{满射}、\Term{open-morphism}{开态射}、\Term{proper}{固有}且\Term{finite-presentation}{有限呈示}，$Y\to S$ 是\Term{finite-presentation}{有限呈示}且\Term{separated}{分离}的。\footnote{编者注：TDTE III，\Term{theorem}{定理} 6.1 指出，若 $S$ 是\Term{noetherian}{诺特}的，则 $Y\to S$ 是\Term{quasi-projective}{拟射影}的。编者未看出如何从\Term{quasi-section}{拟截面}的\Term{local-scheme}{局部}存在性\Term{pushout}{推出}这一点。}

(iii) \Term{set-component}{分量}为 $d_0,d_1$ 的\Term{morphism}{态射} $X_1\to X_0\times_YX_0$ 为\Term{surjectivity}{满射}。

(iv) 若 $(d_0,d_1)$ 是\Term{equivalence-pair}{等价对}，则 $X_1\to X_0\times_YX_0$ 是\Term{isomorphism}{同构}，且 $p$ 是\Term{faithfully-flat}{忠实平坦}的。\footnote{编者注：我们明确写出了后面的推论；参见 [Ray67a]，\Term{theorem}{定理} 1 (iii) 以及\Term{theorem}{定理}证明的末尾。还要指出，对于\Term{equivalence-relation}{等价关系}的情形，[AK80]，\Term{theorem}{定理} 2.9 给出了\Term{theorem}{定理} 7.1 的另一个证明，依据是 Hilbert \Term{scheme}{概形}的存在性；另见 [BLR90]，\S 8.2，\Term{theorem}{定理} 12，其中还证明此时 $Y\to S$ 是\Term{quasi-projective}{拟射影}的。}此外，$(Y,p)$ 是 $(d_0,d_1)$ 在 $(\mathrm{fppf})$ \Term{topology}{拓扑}的\Term{sheaf}{层}\Term{category}{范畴}中的\Term{coequalizer}{余核}；对任意\Term{base-change}{基变换} $Y'\to Y$，$Y'$ 是\Term{groupoid}{群胚} $X_*\times_YY'$ 的\Term{coequalizer}{余核}，后者由 $X_*$ 经\Term{base-change}{基变换} $X_0\times_YY'\to X_0$ 得到。

特别地，对任意\Term{base-change}{基变换} $S'\to S$，$Y'=Y\times_SS'$ 是 $S'$-\Term{groupoid}{群胚} $X'_*=X_*\times_SS'$ 的\Term{coequalizer}{余核}。因此，在这种情形下，“\Term{quotient}{商}的形成与\Term{base-change}{基变换}交换”。
\end{SGAresult}
设 $(Y,p)$ 是 $(d_0,d_1)$ 在全体\Term{ringed-space}{赋环空间}的\Term{category}{范畴}中的\Term{coequalizer}{余核}。\Term{lemma}{引理} 1.2 表明，为证明 (i)，只须证明 $X_0$ 的每个点 $z$ 都有一个\Term{saturated}{饱和}\Term{open-neighborhood}{开邻域} $U_z$，满足如下条件：记 $\widetilde d_0,\widetilde d_1$ 为 $d_0,d_1$ 在 $d_0^{-1}(U_z)=d_1^{-1}(U_z)$ 上的\Term{restriction}{限制}，$(Q,q)$ 为 $(\widetilde d_0,\widetilde d_1)$ 在 $(\mathrm{Esp.An})$ 中的\Term{coequalizer}{余核}，则 $Q$ 是\Term{scheme}{概形}，$q$ 是\Term{scheme-morphism}{概形态射}。

因此，由\Term{lemma}{引理} 6.1 (i)，只须证明 $X_0$ 的每个点 $z$ 都有一个\Term{saturated}{饱和}\Term{open-neighborhood}{开邻域} $U_z$，使 $X_*$ 在 $U_z$ 上诱导的\Term{groupoid}{群胚}有一个\Term{quasi-section}{拟截面}。
\SourcePage{279}
甚至可以假定 $z$ 在其位于 $S$ 上的\Term{fiber}{纤维}中为\Term{closed-point}{闭点}（我们称 $z$ 相对于 $S$ 是闭的）。\footnote{编者注：事实上，若已对每个相对于 $S$ 为闭的点 $z$ 构造了这样的\Term{open-neighborhood}{开邻域} $U_z$，则这些 $U_z$ 的并\Term{covering}{覆盖} $X_0$，因为其闭补集在 $S$ 上的每个\Term{fiber}{纤维}都是没有\Term{closed-point}{闭点}的\Term{noetherian}{诺特}\Term{scheme}{概形}，因而为空。}$U_z$ 的存在性遂由以下\Term{lemma}{引理}得出。

\begin{SGAresult}{\Term{lemma}{引理}}{7.2}
设 $T$ 是\Term{affine}{仿射}\Term{noetherian}{诺特}\Term{scheme}{概形}，$X,Y,Z$ 是\Term{finite-type}{有限型} $T$-\Term{scheme}{概形}，其中 $X$ 在 $T$ 上\Term{quasi-projective}{拟射影}，并设
\[
\begin{tikzcd}
Y\arrow[r,"u"]\arrow[d,"v"']&X\arrow[d,dashed]\\
Z\arrow[r,dashed]&T
\end{tikzcd}
\]
是 $(\mathrm{Sch}/T)$ 中的一个图。另设 $z$ 是 $v(Y)$ 中相对于 $T$ 为闭的一点，且 $v$ 在 $v^{-1}(z)$ 的各点处\Term{flat}{平坦}。则存在 $X$ 的\Term{closed-subscheme}{闭子概形} $F$，使 $u(u^{-1}(F)\cap v^{-1}(z))$ 是\Term{finite}{有限}\Term{nonempty}{非空}集，且 $v$ 在 $u^{-1}(F)$ 上的\Term{restriction}{限制}在 $v^{-1}(z)$ 的各点处\Term{flat}{平坦}。
\end{SGAresult}
\begin{proof}
设 $T=\operatorname{Spec}A$。可以假设 $X$ 形如 $\operatorname{Proj}\mathscr S$，其中 $\mathscr S$ 是一个\Term{finite-type}{有限型} $A$-\Term{module}{模} $E$ 的\Term{symmetric-algebra}{对称代数}。

若 $u(v^{-1}(z))$ \Term{finite}{有限}，可取 $F=X$。否则，以 $y_1,\ldots,y_n$ 表示\Term{fiber}{纤维} $v^{-1}(z)$ 中与其\Term{structure-sheaf}{结构层} $\mathcal O_{v^{-1}(z)}$ \Term{associated-point}{相伴的点}（因此，若 $\mathcal O_i$ 表示 $v^{-1}(z)$ 在 $y_i$ 处的\Term{local-ring}{局部环}，则 $\mathcal O_i$ 的\Term{maximal-ideal}{极大理想}由\Term{zero-divisor}{零因子}组成）。若 $t$ 是 $z$ 在 $T$ 中的像，则 $u(v^{-1}(z))$ 是 $t$ 在 $X$ 中的\Term{fiber}{纤维}的一个无限\Term{constructible-subset}{可构造子集}。故存在这个\Term{fiber}{纤维}中的一个\Term{closed-point}{闭点} $x$，属于 $u(v^{-1}(z))$，且异于 $u(y_1),\ldots,u(y_n)$。那么 $X-\overline{\{x\}}$ 是 $u(y_1),\ldots,u(y_n)$ 的\Term{open-neighborhood}{开邻域}，因而包含一个形如 $D_+(f)$ 的\Term{open-neighborhood}{开邻域}，其中 $f$ 是 $\mathscr S$ 的一个 $d$ 次\Term{homogeneous-element}{齐次元素}（记号同 EGA II，\S 2.3）。

所以，$f$ 定义的\Term{closed-subscheme}{闭子概形} $X_1=V_+(f)$ 包含 $x$，而避开 $u(y_1),\ldots,u(y_n)$。显然，该\Term{subscheme}{子概形}的\Term{inverse-image}{逆像} $Y_1=u^{-1}(V_+(f))$ 异于 $Y$，且与 $v^{-1}(z)$ 相交。我们还将证明 $v$ 在 $Y_1$ 上的\Term{restriction}{限制} $v_1$ 在 $v^{-1}(z)$ 的各点处\Term{flat}{平坦}；于是，若 $u(v_1^{-1}(z))$ \Term{finite}{有限}，只须取 $F=X_1$；否则，将刚才的论证重复一遍，以 $Y_1$ 代替 $Y$，$v_1$ 代替 $v$，以 $u$ 在 $Y_1$ 上诱导的\Term{morphism}{态射} $u_1$ 代替 $u$。这样就得到 $X$ 的一个递降的\Term{closed-subscheme}{闭子概形}序列 $X,X_1,\ldots$；由于这样的序列必定终止，某个 $n$ 将使 $u(u^{-1}(X_n)\cap v^{-1}(z))$ \Term{finite}{有限}\Term{nonempty}{非空}，此时取 $F=X_n$。

因此只须证明 $v_1$ 在 $v^{-1}(z)$ 的各点处\Term{flat}{平坦}。设 $y$ 是 $Y_1$ 中位于 $z$ 上的一点，$\mathcal O_y$ 是 $y$ 在 $Y$ 中的\Term{local-ring}{局部环}，$\overline{\mathcal O}_y$ 是 $y$ 在 $v^{-1}(z)$ 中的\Term{local-ring}{局部环}，$\mathcal O_{v(y)}$ 是 $v(y)$ 在 $Z$ 中的\Term{local-ring}{局部环}。若 $g\in\mathscr S_1$ 使 $D_+(g)$ 是 $u(y)$ 在 $X$ 中的\Term{neighborhood}{邻域}，记 $\varphi$ 为 $f/g^d$ 在 $\mathcal O_y$ 中的像，$\overline\varphi$ 为 $f/g^d$ 在 $\overline{\mathcal O}_y$ 中的像。由 $f$ 的构造可知，$\overline\varphi$ 在 $\overline{\mathcal O}_y$ 中不是\Term{zero-divisor}{零因子}；
\SourcePage{280}
又因 $\mathcal O_y$ 在 $\mathcal O_z$ 上\Term{flat}{平坦}，$\varphi$ 在 $\mathcal O_y$ 中不是\Term{zero-divisor}{零因子}，且 $\mathcal O_y/\mathcal O_y\varphi$ 在 $\mathcal O_z$ 上\Term{flat}{平坦}（SGA 1，IV，5.7）。而 $\mathcal O_y/\mathcal O_y\varphi$ 正是 $y$ 在 $Y_1$ 中的\Term{local-ring}{局部环}。
\end{proof}

\begin{SGAresult}{\Term{lemma}{引理}}{7.3}
沿用 7.1 的记号和假设。则 $X_0$ 的每个相对于 $S$ 为闭的点 $z$ 都有一个\Term{saturated}{饱和}\Term{open-neighborhood}{开邻域} $U_z$，使 $X_*$ 在 $U_z$ 上诱导的\Term{groupoid}{群胚}有一个\Term{quasi-section}{拟截面}。
\end{SGAresult}
\begin{proof}
命题在 $S$ 上是\Term{local-scheme}{局部}的，故可假设 $S$ \Term{affine}{仿射}\Term{noetherian}{诺特}，并将上一\Term{lemma}{引理}应用于 $(\mathrm{Sch}/S)$ 中的图
\[
\begin{tikzcd}
X_1\arrow[r,"d_0"]\arrow[d,"d_1"']&X_0\arrow[d,dashed]\\
X_0\arrow[r,dashed]&S.
\end{tikzcd}
\]
于是，取 $X_0$ 的\Term{closed-subscheme}{闭子概形} $F$，使 $d_0(d_0^{-1}(F)\cap d_1^{-1}(z))$ \Term{finite}{有限}\Term{nonempty}{非空}，且 $d_1$ 在 $d_0^{-1}(F)$ 上的\Term{restriction}{限制}在 $d_1^{-1}(z)$ 的各点处\Term{flat}{平坦}。记 $F_1,F_2$ 为 $F$ 经 $d_0$ 和 $d_0d'_0=d_0d'_1$ 的\Term{inverse-image}{逆像}，并以 $\widetilde d_0,\widetilde d_1$ 等表示 $d_0,d_1$ 等所诱导的\Term{morphism}{态射}。于是有\Term{commutative-diagram}{交换图}
\[
\begin{tikzcd}[column sep=large]
F_2\arrow[r,shift left=.6ex,"\widetilde d'_1"]\arrow[r,shift right=.6ex,"\widetilde d'_0"']\arrow[d,"\widetilde d'_2"']&F_1\arrow[r,"\widetilde d_0"]\arrow[d,"\widetilde d_1"]&F\arrow[d,dashed,"\widetilde q"]\\
X_1\arrow[r,shift left=.6ex,"d_1"]\arrow[r,shift right=.6ex,"d_0"']&X_0\arrow[r,dashed,"q"']&S,
\end{tikzcd}
\]
其中左边两个方块是\Term{cartesian}{笛卡尔}的，第一行是\Term{exact}{正合}的（参见第 1 节的 (0,1,2)），$q,\widetilde q$ 表示\Term{structure-morphism}{结构态射}。

先证明 $F_1$ 在 $z$ 上只有\Term{finite}{有限}个点。\footnote{编者注：这里增添了一些细节，并突出了\Term{theorem}{定理} 7.1 中 $d_0,d_1$ 的\Term{proper}{固有}性假设所起的\Term{action}{作用}。（可以与\Term{theorem}{定理} 8.1 的陈述和证明比较，后者省去了这个\Term{proper}{固有}性假设。）}事实上，设 $s$ 为 $z$ 在 $S$ 中的像；因为 $F$ 在 $S$ 上\Term{finite-type}{有限型}，\Term{fiber}{纤维} $\widetilde q^{-1}(s)$ 是\Term{noetherian}{诺特}\Term{scheme}{概形}。另一方面，因为 $\widetilde d_0$ 是\Term{proper}{固有}的，$\widetilde d_0(\widetilde d_1^{-1}(z))$ 是 $\widetilde q^{-1}(s)$ 的一个\Term{closed-subscheme}{闭子概形}，由\Term{finite}{有限}个点组成。因此（参见 EGA I，6.2.2），这个\Term{set}{集合}的各点在 $\widetilde q^{-1}(s)$ 中为\Term{closed-point}{闭点}，而且（因为 $F$ 在 $X_0$ 中闭）在 $s$ 于 $X_0$ 中的\Term{fiber}{纤维} $q^{-1}(s)$ 中也是\Term{closed-point}{闭点}。设 $y$ 为这些点之一；因为\Term{fiber}{纤维} $q^{-1}(s)$ 在 $\kappa(s)$ 上\Term{finite-type}{有限型}，其中分别存在 $y,z$ 的\Term{affine-open}{仿射开邻域} $\operatorname{Spec}B,\operatorname{Spec}C$，$B,C$ 为\Term{finite-type}{有限型} $\kappa(s)$-\Term{algebra}{代数}。于是 $y,z$ 分别对应\Term{maximal-ideal}{极大理想} $\mathfrak p\subset B,\mathfrak q\subset C$，\Term{field}{域} $B/\mathfrak p,C/\mathfrak q$ 在 $\kappa(s)$ 上的\Term{degree}{次数}\Term{finite}{有限}，所以 $(B/\mathfrak p)\otimes_{\kappa(s)}(C/\mathfrak q)$ 是\Term{finite-dimensional}{有限维} $\kappa(s)$-\Term{algebra}{代数}，其\Term{maximal-ideal}{极大理想}恰好对应于 $X_0\times_SX_0$ 中第二（分别第一）\Term{projection}{投影}为 $z$（分别为 $y$）的点。因此，$X_0\times_SX_0$ 中第二\Term{projection}{投影}为 $z$、第一\Term{projection}{投影}属于 $\widetilde d_0(\widetilde d_1^{-1}(z))$ 的点 $u$ 只有\Term{finite}{有限}个。最后，因为 $X_1\to X_0\times_SX_0$ 的\Term{fiber}{纤维}\Term{finite}{有限}，
\SourcePage{281}
这样的点 $u$ 只来自 $X_1$ 的\Term{finite}{有限}个点，所需断言遂得证。

因此，\Term{morphism}{态射} $\widetilde d_1$ 在 $F_1$ 位于 $z$ 上的各点处\Term{quasi-finite}{拟有限}且\Term{flat}{平坦}。因为 $\widetilde d_1$ 为\Term{finite-type}{有限型}，由 SGA 1，IV，6.10 及 EGA III，4.4.10，\footnote{编者注：若 $f:X\to Y$ 为\Term{locally-finite-type}{局部有限型}\Term{morphism}{态射}，则 $X$ 中在其\Term{fiber}{纤维} $f^{-1}(f(x))$ 中孤立的点 $x$ 所成的\Term{set}{集合}在 $X$ 中开：当 $Y$ 为\Term{locally-noetherian}{局部诺特}时，EGA III，4.4.10 从 Zariski 的“主\Term{theorem}{定理}”\Term{pushout}{推出}这一点；对任意 $Y$，这由 Chevalley 的\Term{semicontinuity-theorem}{半连续性定理}得到（EGA IV$_3$，13.1.3 和 13.1.4）。因此，$f$ 在 $x$ 处\Term{quasi-finite}{拟有限}，当且仅当 $f$ 在 $x$ 处\Term{finite-type}{有限型}且 $x$ 在 $f^{-1}(f(x))$ 中孤立；这一点以后还将用到，参见编者注 (42)。}$F_1$ 中 $\widetilde d_1$ 不同时\Term{flat}{平坦}且\Term{quasi-finite}{拟有限}的点所成的\Term{set}{集合} $\Phi$ 在 $F_1$ 中闭，因而在 $X_1$ 中闭（因为 $F_1$ 在 $X_1$ 中闭）。因为 $d_1$ 是\Term{proper}{固有}的，$\widetilde d_1(\Phi)$ 为闭集，且由前述结果不包含 $z$。置 $W=\widetilde d_1(F_1)-\widetilde d_1(\Phi)$。那么，$\widetilde d_1$ 在 $\widetilde d_1^{-1}(W)$ 上的\Term{restriction}{限制}\footnote{编者注：以下展开了原文的细节。}是\Term{finite-presentation}{有限呈示}的（由\Term{noetherian}{诺特}性假设）、\Term{flat}{平坦}、\Term{proper}{固有}且\Term{quasi-finite}{拟有限}，因而由 EGA III，4.4.2，以及 EGA IV$_2$，2.1.12 和 2.4.6，它是\Term{finite-morphism}{有限}、\Term{locally-free}{局部自由}且为\Term{open-morphism}{开态射}。故 $\widetilde d_1(F_1)$ 是 $z$ 的一个\Term{neighborhood}{邻域}，而 $W$ 是 $X_0$ 中包含在 $\widetilde d_1(F_1)$ 内、且在其上方 $\widetilde d_1$ 同时\Term{flat}{平坦}和\Term{quasi-finite}{拟有限}的最大\Term{open-subset}{开子集}。

我们将在\Term{lemma}{引理} 7.4 中看到，$\Phi$ 经 $\widetilde d'_1,\widetilde d'_0$ 的\Term{inverse-image}{逆像}都等同于 $F_2$ 中 $\widetilde d'_2$ 不同时\Term{flat}{平坦}和\Term{quasi-finite}{拟有限}的点所成的\Term{set}{集合}。因此
\[
d_0^{-1}(W)=\widetilde d'_2(F_2)-\widetilde d'_2(\widetilde d_0'{}^{-1}\Phi)
\]
与
\[
d_1^{-1}(W)=\widetilde d'_2(F_2)-\widetilde d'_2(\widetilde d_1'{}^{-1}\Phi)
\]
相同，即 $W$ 是\Term{saturated}{饱和}的。故置 $W_1=\widetilde d_1^{-1}(W)$，等式 $d_0^{-1}(W)=d_1^{-1}(W)$ 给出
\[
\widetilde d_2'{}^{-1}d_0^{-1}(W)=\widetilde d_2'{}^{-1}d_1^{-1}(W),
\quad\text{即}\quad
\widetilde d_0'{}^{-1}(W_1)=\widetilde d_1'{}^{-1}(W_1).
\]
因为 $\widetilde d_0$ 是\Term{faithfully-flat}{忠实平坦}且\Term{quasi-compact}{拟紧}的（因为 $d_0$ 与 $d_1$ 一样，为\Term{surjectivity}{满射}、\Term{proper}{固有}且\Term{flat}{平坦}），并且方块
\[
\begin{tikzcd}
F_2\arrow[r,"\widetilde d'_1"]\arrow[d,"\widetilde d'_0"']&F_1\arrow[d,"\widetilde d_0"]\\
F_1\arrow[r,"\widetilde d_0"']&F
\end{tikzcd}
\]
是\Term{cartesian}{笛卡尔}的，故 $W_1$ 形如 $\widetilde d_0^{-1}(U)$，其中 $U$ 是 $F$ 的一个\Term{open-subset}{开子集}（SGA 1，VIII，4.4）。$F$ 的这个\Term{open-subset}{开子集} $U$ 是 $X_*$ 所诱导的以 $W$ 为底的\Term{groupoid}{群胚}的\Term{quasi-section}{拟截面}。因此可取 $U_z=W$。
\end{proof}
现在只须陈述\Term{lemma}{引理} 7.4，其证明是经典的。
\begin{SGAresult}{\Term{lemma}{引理}}{7.4}
考虑\Term{scheme}{概形}的一个\Term{cartesian-square}{笛卡尔方块}
\[
\begin{tikzcd}
F_2\arrow[r,"v"]\arrow[d,"d'"']&F_1\arrow[d,"d"]\\
X_1\arrow[r,"u"']&X_0
\end{tikzcd}
\]
\SourcePage{282}
并设 $x$ 是 $F_2$ 的一点。

(i) 若 $u$ \Term{flat}{平坦}，则 $d'$ 在 $x$ 处\Term{flat}{平坦}，当且仅当 $d$ 在 $v(x)$ 处\Term{flat}{平坦}。

(ii) 若 $d$ 是\Term{locally-finite-type}{局部有限型}的，则 $d'$ 在 $x$ 处\Term{quasi-finite}{拟有限}，当且仅当 $d$ 在 $v(x)$ 处\Term{quasi-finite}{拟有限}。\footnote{编者注：由\Term{base-change}{基变换}可知这些条件是充分的（参见 EGA II，6.2.4 (iii) 和 EGA IV$_2$，2.1.4）。反过来，置 $y=d'(x)$、$z=u(y)=d(v(x))$，并假设 $d'$ 在 $x$ 处\Term{flat}{平坦}，$u$（故 $v$ 也）\Term{flat}{平坦}。于是 $\mathcal O_{v(x)}\to\mathcal O_x$ \Term{faithfully-flat}{忠实平坦}，$\mathcal O_z\to\mathcal O_y\to\mathcal O_x$ 也\Term{faithfully-flat}{忠实平坦}。所以 $\mathcal O_z\to\mathcal O_{v(x)}$ \Term{faithfully-flat}{忠实平坦}（参见 EGA IV$_2$，2.2.11 (iv)）。最后，假设 $d$ 为\Term{locally-finite-type}{局部有限型}且 $d'$ 在 $x$ 处\Term{quasi-finite}{拟有限}。因为 $x$ 在其\Term{fiber}{纤维} $d'{}^{-1}(y)=d^{-1}(z)\otimes_{\kappa(z)}\kappa(y)$ 中孤立，$v(x)$ 在其\Term{fiber}{纤维} $d^{-1}(z)$ 中也孤立。因此，由 Chevalley 的\Term{semicontinuity-theorem}{半连续性定理}，存在 $v(x)$ 的一个\Term{open-neighborhood}{开邻域}，其中每一点都在其\Term{fiber}{纤维}中孤立（EGA IV$_3$，13.1.3 和 13.1.4），故 $d$ 在 $v(x)$ 处\Term{quasi-finite}{拟有限}。}
\end{SGAresult}
我们已证明，$X_0$ 可由\Term{saturated}{饱和}\Term{open-subset}{开子集} $W$ \Term{covering}{覆盖}，使 $X_*$ 在 $W$ 上诱导的\Term{groupoid}{群胚} $W_*$ 有\Term{quasi-section}{拟截面}。\footnote{编者注：以下利用\Term{lemma}{引理} 6.1 中的增补修改了原文。}

由\Term{lemma}{引理} 6.1 以及\Term{theorem}{定理} 7.1 后所述的化简，可得\Term{theorem}{定理} 7.1 的断言 (i)、(iii)，并且 $p$ 为\Term{surjectivity}{满射}和\Term{open-morphism}{开态射}，$p$ 与 $Y\to S$ 都是\Term{locally-finite-presentation}{局部有限呈示}的。此外，因 $X_0\to S$ \Term{quasi-projective}{拟射影}，故\Term{separated}{分离}且\Term{finite-type}{有限型}，于是 $p$ \Term{separated}{分离}，而\Term{lemma}{引理} 6.1 (ii) 的证明表明，$p$ 与 $Y\to S$ 是\Term{finite-presentation}{有限呈示}的。

为证明 $p$ \Term{proper}{固有}，只须再证明它\Term{universally-closed}{普遍闭}。这个断言在 $Y$ 上是\Term{local-scheme}{局部}的，故可在一个\Term{saturated}{饱和}\Term{open-subset}{开子集} $W$ 上工作，使 $X_*$ 在 $W$ 上诱导的\Term{groupoid}{群胚} $W_*$ 有\Term{quasi-section}{拟截面} $U$（因为这样的\Term{open-subset}{开子集}\Term{covering}{覆盖} $X_0$）。沿用第 6 节 a) 的记号，有\Term{commutative-diagram}{交换图}
\[
\begin{tikzcd}
W\arrow[dr,"p"']&V\arrow[l,"v"']\arrow[r,"u"]\arrow[d,"r"]&U\arrow[dl,"q"]\\
&Z,&
\end{tikzcd}
\]
其中 $Z$ 是 $Y$ 的\Term{open-subset}{开子集}，所有\Term{arrow}{箭头}均为\Term{surjectivity}{满射}，$q$ 是\Term{integral-morphism}{整态射}。又由假设 $d_0:X_1\to X_0$ \Term{proper}{固有}，所以经\Term{base-change}{基变换}由它得到的 $u$ 也\Term{proper}{固有}。故 $r$ \Term{universally-closed}{普遍闭}；因 $v$ 为\Term{surjectivity}{满射}，$p$ 也\Term{universally-closed}{普遍闭}。

最后，$p$ 为\Term{surjectivity}{满射}且\Term{universally-closed}{普遍闭}，$X_0$ \Term{quasi-projective}{拟射影}因而\Term{separated}{分离}，故\Term{diagonal}{对角线} $\Delta_{Y/S}(Y)$ 在 $Y\times_SY$ 中闭，因为它是\Term{diagonal}{对角线} $\Delta_{X_0/S}(X_0)$ 经 $p\times p$ 的像。所以 $Y$ 在 $S$ 上\Term{separated}{分离}。\Term{theorem}{定理} 7.1 (ii) 得证。

7.1 (iv) 中待证的断言在 $Y$ 上是\Term{local-scheme}{局部}的；$X_0$ 由\Term{saturated}{饱和}\Term{open-subset}{开子集} $U_z$ \Term{covering}{覆盖}，故只须以 $U_z$ 和 $V=p(U_z)$ 代替 $X$ 和 $Y$ 来证明这些断言。如在 7.1 的证明开头已经看到的，由\Term{lemma}{引理} 1.1、1.2 和 6.1 (i)，$(V_z,p|_{U_z})$ 是 $X_*$ 在 $U_z$ 上诱导的\Term{groupoid}{群胚}在 $(\mathrm{Sch})$ 及 $(\mathrm{Esp.An})$ 中的\Term{coequalizer}{余核}。而 $d=(d_0,d_1)$ 是\Term{monomorphism}{单态射}这一假设在\Term{base-change}{基变换} $U_z\to X_0$ 下保持。
\SourcePage{283}
所以，7.1 (iv) 的前两个断言由 6.1 (iv) 得出。

最后证明 (iv) 末尾所指出的推论（参见 [Ray67a]，\Term{theorem}{定理} 1 (iii)）。由假设，\Term{groupoid}{群胚} $X_*$ 来自\Term{equivalence-relation}{等价关系} $R\to X_0\times X_0$；我们已证明 $R$ 是\Term{effective-equivalence-relation}{有效}的（参见第 IV 讲 3.3.2），且 $p:X_0\to Y=X_0/R$ \Term{faithfully-flat}{忠实平坦}且\Term{finite-presentation}{有限呈示}。因此，以 $(M)$ 表示\Term{faithfully-flat}{忠实平坦}\Term{locally-finite-presentation}{局部有限呈示}\Term{morphism}{态射}族，则 $R$ 是 $(M)$-\Term{effective-equivalence-relation}{有效}的。所以由第 IV 讲 6.3.3，$(Y,p)$ 表示 $X_0$ 对 $R$ 的 $(\mathrm{fppf})$ \Term{topology}{拓扑}的\Term{quotient-sheaf}{商层}，而关于\Term{base-change}{基变换}的断言由 IV，3.4.3.1 得出。

\SGAsection{8}{对未必固有的平坦群胚取商}
\begin{SGAresult}{\Term{theorem}{定理}}{8.1}
\footnote{编者注：$X_0$ 中存在满足\Term{theorem}{定理}结论的最大\Term{open-subset}{开子集} $W$。事实上，设 $W$ 是\Term{theorem}{定理}中的\Term{open-subset}{开子集}，$W^\sharp$ 是包含于 $W$ 的\Term{dense}{稠密}\Term{saturated}{饱和}\Term{open-subset}{开子集}。因为 $p$ 为\Term{open-morphism}{开态射}，$V^\sharp=p(W^\sharp)$ 是 $V$ 的\Term{open-subset}{开子集}；因为 $W^\sharp$ \Term{saturated}{饱和}，有 $W^\sharp=p^{-1}(V^\sharp)$。由\Term{lemma}{引理} 1.1，$V^\sharp$ 是 $W^\sharp$ 上诱导的\Term{groupoid}{群胚}的\Term{coequalizer}{余核}。于是可沿交集 $W^\sharp$ \Term{gluing}{粘合}两个满足\Term{theorem}{定理}结论的\Term{open-subset}{开子集} $W,W'$，条件 (i)、(ii)、(iii)、(iv)，以及 $p$ 和 $V\to S$ 为\Term{locally-finite-presentation}{局部有限呈示}这一性质都得到保持。结论 (ii') 如 6.1 (ii) 的证明一样，由 $X_0$ 在\Term{noetherian}{诺特}\Term{scheme}{概形} $S$ 上\Term{finite-type}{有限型}这一假设得出。
另一方面，\Term{lemma}{引理} 1.1 和 1.2 还表明，$X_0$ 的所有满足 $p(U)$ 是\Term{scheme}{概形}且 $p|_U:U\to p(U)$ 是\Term{scheme-morphism}{概形态射}的\Term{saturated}{饱和}\Term{open-subset}{开子集} $U$ 的并，是 $X_0$ 中满足 8.1 条件 (i) 的最大\Term{saturated}{饱和}\Term{open-subset}{开子集} $\Omega$。\Term{theorem}{定理} 8.1 表明 $\Omega$ 包含一个\Term{dense}{稠密}\Term{open-subset}{开子集} $W$，但 $\Omega$ 是否满足 (ii) 至 (iv) 并不立即清楚。
关于这一点，读者可参阅 [Ray67a]、[Ray67b] 和 [An73] 的附录 I，它们给出了更精确的结果，并研究 $S$-\Term{quotient-sheaf}{商层} $\widetilde{X/R}$ 的 $(\mathrm{fppf})$ \Term{representability}{可表性}（这里以 $R$ 表示以 $X=X_0$ 为底的\Term{groupoid}{群胚} $X_*$）；所用假设更弱：$S$ 是任意\Term{scheme}{概形}，$X$ 是在 $S$ 上\Term{locally-finite-type}{局部有限型}的\Term{scheme}{概形}，$R$ 是以 $X$ 为底的 $S$-\Term{groupoid}{群胚}，$d_0$（因而 $d_1$）\Term{flat}{平坦}且\Term{finite-presentation}{有限呈示}。特别指出下列结果：若 $\widetilde{X/R}$ 可由一个 $S$-\Term{scheme}{概形} $Y$ 表示，则 $Y$ 也是 $(\mathrm{Esp.An})$ 中的\Term{coequalizer}{余核}。反命题一般不成立（参见 [Mum65]，第 0 章 \S 3 的例 0.4，[Ray67a]，注 1 引用了此例），但若 $d=(d_0,d_1)$ 是\Term{immersion}{浸入}，反命题成立。在此假设下，\Term{morphism}{态射} $p:\Omega\to Z:=\Omega/R_\Omega$ \Term{faithfully-flat}{忠实平坦}且\Term{finite-presentation}{有限呈示}；若再假设 $S$ \Term{locally-noetherian}{局部诺特}，则 $X$ 中\Term{codimension}{余维}为 $1$ 的点 $x$ 属于 $\Omega$，当且仅当在 $\operatorname{Spec}(\mathcal O_{X,x})$ 上诱导的\Term{groupoid}{群胚}的\Term{graph}{图像}是闭的。以上各点参见 [Ray67a]，命题 1；[Ray67b]，命题 1 及\Term{theorem}{定理} 2、1、4；[An73]，第 66--67 页\Term{theorem}{定理} 5、6 及第 49 页命题 3.3.1。（对于在\Term{algebraically-closed-field}{代数闭域} $k$ 上的\Term{algebraic-group}{代数群}\Term{action}{作用}，另见文章 [DR81]。）}
设 $S$ 是\Term{noetherian}{诺特}\Term{scheme}{概形}，且
\[
\begin{tikzcd}
X_2\arrow[r,shift left=1.1ex,"d'_2"]\arrow[r,"d'_1" description]\arrow[r,shift right=1.1ex,"d'_0"']&X_1\arrow[r,shift left=.6ex,"d_1"]\arrow[r,shift right=.6ex,"d_0"']&X_0
\end{tikzcd}
\]
是 $(\mathrm{Sch}/S)$-\Term{groupoid}{群胚}，其中 $d_1$ \Term{flat}{平坦}且\Term{finite-type}{有限型}，$X_0$ 在 $S$ 上\Term{finite-type}{有限型}，\Term{set-component}{分量}为 $d_0,d_1$ 的\Term{morphism}{态射} $X_1\to X_0\times_SX_0$ \Term{quasi-finite}{拟有限}。则 $X_0$ 中存在\Term{dense}{稠密}\Term{saturated}{饱和}\Term{open-subset}{开子集} $W$，满足下列性质：

(i) 若 $W_2\rightrightarrows W_1\rightrightarrows W$ 是 $X_*$ 在 $W$ 上诱导的\Term{groupoid}{群胚}，\Term{arrow}{箭头}分别记作 $w'_i,w_j$，则 $(w_0,w_1)$ 在 $(\mathrm{Sch}/S)$ 中有\Term{coequalizer}{余核} $(V,p)$；而且，$(V,p)$ 也是 $(w_0,w_1)$ 在全体\Term{ringed-space}{赋环空间}的\Term{category}{范畴}中的\Term{coequalizer}{余核}。

(ii) $p$ 为\Term{surjectivity}{满射}和\Term{open-morphism}{开态射}。

(ii') $p$ 与 $V\to S$ 为\Term{finite-presentation}{有限呈示}。

(iii) \Term{set-component}{分量}为 $w_0,w_1$ 的\Term{morphism}{态射} $W_1\to W\times_VW$ 为\Term{surjectivity}{满射}。

(iv) 若 $(d_0,d_1)$ 是\Term{equivalence-pair}{等价对}，则 $W_1\to W\times_VW$ 是\Term{isomorphism}{同构}，且 $p$ \Term{faithfully-flat}{忠实平坦}。
\end{SGAresult}
\SourcePage{284}
\begin{proof}
我们将证明，可以选择 $W$，使 $X_*$ 在其上诱导的 $(\mathrm{Sch}/S)$-\Term{groupoid}{群胚} $W_*$ 有\Term{quasi-section}{拟截面}（参见第 7 节）。\Term{theorem}{定理} 8.1 随即由\Term{lemma}{引理} 6.1 得出。

暂且承认，对 $X_0$ 中每个相对于 $S$ 为闭的点 $z$（参见第 7 节），存在\Term{saturated}{饱和}\Term{open-subset}{开子集} $W_z$，它有\Term{quasi-section}{拟截面}，并与 $X_0$ 中经过 $z$ 的每个\Term{irreducible-component}{不可约分支}相交。那么，$W_z$ 在 $X_0$ 中的外部 $X_0-\overline{W_z}$ 是\Term{saturated}{饱和}的（因为此外部的\Term{saturation}{饱和化} $d_1(d_0^{-1}(X_0-\overline{W_z}))$ 是开的，且不与 $W_z$ 相交）。若此外部\Term{nonempty}{非空}，可在其中选择相对于 $S$ 为闭的点 $z'$，并对 $z'$ 取如上所述的\Term{open-subset}{开子集} $W_{z'}$；而且可假设 $W_{z'}$ 包含于 $X_0-\overline{W_z}$。则 $W_z$ 与 $W_{z'}$ 不交，$X_*$ 在 $W_z\cup W_{z'}$ 上诱导的\Term{groupoid}{群胚}有\Term{quasi-section}{拟截面}。这个过程必定终止，因为 $X_0$ 只有\Term{finite}{有限}个\Term{irreducible-component}{不可约分支}。因此只须构造 $W_z$。

为此可假设 $S$ \Term{affine}{仿射}；取 $X_1$ 的一点 $y$ 使 $d_1(y)=z$，取 $X_0$ 中包含 $d_0(y)$ 的一个\Term{affine-open}{仿射开子集} $X$，记 $Y$ 为 $X$ 经 $d_0$ 在 $X_1$ 中的\Term{inverse-image}{逆像}，$u:Y\to X$、$v:Y\to X_0$ 分别为 $d_0,d_1$ 诱导的\Term{morphism}{态射}。因为 $X$ \Term{affine}{仿射}，故\Term{quasi-projective}{拟射影}，可应用\Term{lemma}{引理} 7.2：因此存在 $X_0$ 的\Term{subscheme}{子概形} $F$，使 $d_0^{-1}(F)\cap d_1^{-1}(z)$ \Term{nonempty}{非空}，$d_0(d_0^{-1}(F)\cap d_1^{-1}(z))$ \Term{finite}{有限}，且 $d_1$ 在 $d_0^{-1}(F)$ 上的\Term{restriction}{限制}在 $v^{-1}(z)$ 的各点处\Term{flat}{平坦}。于是可沿用\Term{lemma}{引理} 7.3 的记号，以 $F_1,F_2$ 表示 $F$ 在 $X_1,X_2$ 中的\Term{inverse-image}{逆像}，等等。
\[
\begin{tikzcd}
F_2\arrow[r,shift left=.6ex,"\widetilde d'_1"]\arrow[r,shift right=.6ex,"\widetilde d'_0"']\arrow[d,"\widetilde d'_2"']&F_1\arrow[r,"\widetilde d_0"]\arrow[d,"\widetilde d_1"]&F\\
X_1\arrow[r,shift left=.6ex,"d_1"]\arrow[r,shift right=.6ex,"d_0"']&X_0.&
\end{tikzcd}
\]
与 7.3 一样，可证明 $\widetilde d_1$ 在 $\widetilde d_1^{-1}(z)$ 的各点处\Term{quasi-finite}{拟有限}，因此自然考虑 $F_1$ 中由 $\widetilde d_1$ 同时\Term{flat}{平坦}且\Term{quasi-finite}{拟有限}的各点组成的\Term{open-subset}{开子集} $F'_1$。由 7.4，$F'_1$ 经 $\widetilde d'_1$ 和 $\widetilde d'_0$ 的两个\Term{inverse-image}{逆像}都由 $F_2$ 中 $\widetilde d'_2$ \Term{flat}{平坦}且\Term{quasi-finite}{拟有限}的点组成，
\SourcePage{285}
故这两个\Term{inverse-image}{逆像}相同，$F'_1$ 形如 $\widetilde d_0^{-1}(F')$，其中 $F'$ 是 $F$ 的一个\Term{open-subset}{开子集}（SGA 1，VIII，4.4）。以 $F'$ 代替 $F$，即可假设 $\widetilde d_1$ \Term{quasi-finite}{拟有限}且\Term{flat}{平坦}。此时，以 $W_z$ 表示 $\widetilde d_1(F_1)$ 中使 $\widetilde d_1$ 在其上方\Term{finite-morphism}{有限}且\Term{flat}{平坦}的最大\Term{open-subset}{开子集}。

这个\Term{open-subset}{开子集} $W_z$ 未必包含 $z$，但包含 $X_0$ 中经过 $z$ 的各\Term{irreducible-component}{不可约分支}的\Term{generic-point}{泛点}。\footnote{编者注：事实上，设 $\eta$ 是这样的\Term{generic-point}{泛点}。由假设，$\mathcal O_{X_0,\eta}$ 是\Term{artinian-local-ring}{阿廷局部环}，$\mathcal O_{F_1,\eta}$ 是\Term{finite-type}{有限型} $\mathcal O_{X_0,\eta}$-\Term{module}{模}。故由 SGA 1，VIII，6.5，存在 $\eta$ 的一个\Term{open-neighborhood}{开邻域}，使 $\widetilde d_1$ 在其上方\Term{finite-morphism}{有限}。}因为 $d_0$（分别 $d_1$）\Term{faithfully-flat}{忠实平坦}且\Term{finite-presentation}{有限呈示}（因而为\Term{open-morphism}{开态射}），由 SGA 1，VIII，5.7，$d_0^{-1}(W_z)$ 和 $d_1^{-1}(W_z)$ 都等于 $\widetilde d'_2(F_2)$ 中使 $\widetilde d'_2$ 在其上方\Term{finite-morphism}{有限}且\Term{flat}{平坦}的最大\Term{open-subset}{开子集}。因此，同 7.3 可知，$\widetilde d_1^{-1}(W_z)$ 经 $\widetilde d'_0,\widetilde d'_1$ 的两个\Term{inverse-image}{逆像}相同，故 $\widetilde d_1^{-1}(W_z)$ 形如 $\widetilde d_0^{-1}(U)$，其中 $U$ 是 $F$ 的\Term{open-subset}{开子集}，并且是 $X_*$ 在 $W_z$ 上诱导的\Term{groupoid}{群胚}的\Term{quasi-section}{拟截面}。
\end{proof}

\SGAsection{9}{消去定理 7.1 中的诺特性假设}
a) 沿用\Term{lemma}{引理} 6.1 的记号和假设，并设 $\pi:S'\to S$ 是任意\Term{base-change}{基变换}。对于 $S$-\Term{scheme-morphism}{概形态射} $f:X\to Y$，记 $f':X'\to Y'$ 为经底的\Term{group-extension}{扩张} $\pi$ 得到的 $S'$-\Term{scheme-morphism}{概形态射}。按此约定，$p':X'_0\to Y'$ 为\Term{surjectivity}{满射}，\Term{set-component}{分量}为 $d'_0,d'_1$ 的\Term{morphism}{态射} $X'_1\to X'_0\times_{Y'}X'_0$ 也为\Term{surjectivity}{满射}。因此，$Y'$ 的底层\Term{set}{集合}可与 $X'_0$ 的底层\Term{set}{集合}对 $S'$-\Term{groupoid}{群胚} $X'_*$ 在 $X'_0$ 中定义的\Term{equivalence-relation}{等价关系}的\Term{quotient}{商}等同。此外，$q':U'\to Y'$ 是\Term{surjectivity}{满射}\Term{integral-morphism}{整态射}，故 $Y'$ 的\Term{topology}{拓扑}是 $U'$ 的\Term{quotient-topology}{商拓扑}，因而也是 $X'_0$ 的\Term{quotient-topology}{商拓扑}（参见第 6 节 c) 的证明）。

另一方面，显然 $U'$ 是 $S'$-\Term{groupoid}{群胚} $X'_*$ 的\Term{quasi-section}{拟截面}，故可对其应用\Term{lemma}{引理} 6.1。特别地，$X'_*$ 有\Term{coequalizer}{余核} $(Y_1,p_1)$，而 $Y_1$ 的底层\Term{topological-space}{拓扑空间}由 $X'_0$ 的底层\Term{topological-space}{拓扑空间}对 $X'_*$ 定义的关系等同\Term{equivalence}{等价}点而得到。因此，\Term{canonical}{典范}\Term{morphism}{态射} $Y_1\to Y'$ 是\Term{homeomorphism}{同胚}；我们进一步断言，$Y_1\to Y'$ 是\Term{universal-homeomorphism}{普遍同胚}：事实上，若 $S''$ 在 $S'$ 上，设 $Y_2$ 是 $(d_0\times_SS'',d_1\times_SS'')$ 的\Term{coequalizer}{余核}。将前述结果应用于\Term{base-change}{基变换} $S''\to S'$ 和 $S''\to S$，可知
\[
Y_2\longrightarrow Y_1\times_{S'}S''
\quad\text{及}\quad
Y_2\longrightarrow Y\times_SS''\simeq Y'\times_{S'}S''
\]
都是\Term{homeomorphism}{同胚}，故 $Y_1\times_{S'}S''\to Y'\times_{S'}S''$ 也是\Term{homeomorphism}{同胚}。

b) 类似的论述显然也适用于\Term{groupoid}{群胚} $X_*$ “\Term{local-scheme}{局部}”具有\Term{quasi-section}{拟截面}的情形（参见\Term{theorem}{定理} 7.1 的证明）。\footnote{编者注：以下展开了细节，以突出下面的\Term{theorem}{定理} 9.0。}例如，有如下\Term{theorem}{定理}。
\SourcePage{286}
\begin{SGAresult}{\Term{theorem}{定理}}{9.0}
设 $S$ 是任意\Term{scheme}{概形}，$X_2\xrightarrow{d'_j}X_1\xrightarrow{d_i}X_0$ 是一个 $(\mathrm{Sch}/S)$-\Term{groupoid}{群胚}，满足：$X_0$ 在 $S$ 上\Term{finite-presentation}{有限呈示}且\Term{quasi-projective}{拟射影}，$d_1$ \Term{finite-presentation}{有限呈示}、\Term{proper}{固有}且\Term{flat}{平坦}，\Term{morphism}{态射} $d_0\boxtimes d_1:X_1\to X_0\times_SX_0$ \Term{quasi-finite}{拟有限}。那么：

(1) $X_0$ 的每个点 $x$ 都有一个\Term{saturated}{饱和}\Term{open-neighborhood}{开邻域} $W$，使 $X_*$ 在 $W$ 上诱导的\Term{groupoid}{群胚}具有\Term{quasi-section}{拟截面}。

(2) 设 $(Y,p)$ 是 $(d_0,d_1)$ 在全体\Term{ringed-space}{赋环空间}的\Term{category}{范畴}中的\Term{coequalizer}{余核}。则 $Y$ 是\Term{scheme}{概形}，$p$ 是\Term{scheme-morphism}{概形态射}，且 $(Y,p)$ 是 $(d_0,d_1)$ 在 $(\mathrm{Sch}/S)$ 中的\Term{coequalizer}{余核}。

(3) $p$ 为\Term{surjectivity}{满射}、\Term{open-morphism}{开态射}且\Term{universally-closed}{普遍闭}。

(4) \Term{set-component}{分量}为 $d_0,d_1$ 的\Term{morphism}{态射} $d:X_1\to X_0\times_YX_0$ 为\Term{surjectivity}{满射}。

(5) 若 $(d_0,d_1)$ 是\Term{equivalence-pair}{等价对}，则：

(a) $d:X_1\to X_0\times_YX_0$ 是\Term{isomorphism}{同构}，$p$ \Term{faithfully-flat}{忠实平坦}。

(b) $p$ 和 $Y\to S$ 都\Term{finite-presentation}{有限呈示}，且 $(Y,p)$ 是 $(d_0,d_1)$ 在 $(\mathrm{fppf})$ \Term{topology}{拓扑}的\Term{sheaf}{层}\Term{category}{范畴}中的\Term{coequalizer}{余核}。
\end{SGAresult}
\begin{proof}
对于 (1)，问题在 $S$ 上是\Term{local-scheme}{局部}的，故可假设 $S=\operatorname{Spec}B$ \Term{affine}{仿射}。于是存在一个在 $\mathbb Z$ 上\Term{finite-type}{有限型}的\Term{ring}{环} $A$、一个\Term{morphism}{态射} $S\to T=\operatorname{Spec}A$ 及一个 $(\mathrm{Sch}/T)$-\Term{groupoid}{群胚} $Z_*$，使 $X_*$ 与 $Z_*\times_TS$ 等同（参见 EGA IV$_3$，8.8.3，应用于 $S_0=\operatorname{Spec}\mathbb Z$ 及 $S_i=\operatorname{Spec}A_i$，其中 $A_i$ 遍历 $B$ 的各\Term{finite-type}{有限型} $\mathbb Z$-\Term{subalgebra}{子代数}）。而且，可假设 $Z_*$ 满足\Term{theorem}{定理} 7.1 的条件（参见 EGA IV$_3$，8.10.5）。故 $Z_*$ “\Term{local-scheme}{局部}”具有\Term{quasi-section}{拟截面}。

由 a)，$X_*$ 也如此，断言 (2)、(3)、(4)、(5)(a) 同\Term{theorem}{定理} 7.1 的证明一样，由 6.1 得出。

c) 证明 $Y\to S$ \Term{finite-presentation}{有限呈示}。\footnote{编者注：原文说这由后面的命题 9.1 得出。我们未能重建这一论证。以下证明由 O. Gabber 告知我们。}由假设，$(d_0^{X_*},d_1^{X_*})$ 是\Term{equivalence-pair}{等价对}，即 $d^{X_*}:X_1\to X_0\times_SX_0$ 是\Term{monomorphism}{单态射}。由 EGA IV$_3$，8.10.5，将 $A$ 适当扩大后，可假设 $d^{Z_*}:Z_1\to Z_0\times_TZ_0$ 是\Term{monomorphism}{单态射}。因为 $T=\operatorname{Spec}A$ 且 $A$ \Term{noetherian}{诺特}，\Term{theorem}{定理} 7.1 表明，\Term{groupoid}{群胚} $Z_*$ 在 $(\mathrm{Sch}/T)$ 中有\Term{coequalizer}{余核} $(Q,q)$，$q$ 与 $Q\to T$ \Term{finite-presentation}{有限呈示}，并且 $q:Z_0\to Q$ \Term{faithfully-flat}{忠实平坦}，$d^{Z_*}$ 诱导\Term{isomorphism}{同构} $Z_1\xrightarrow\sim Z_0\times_QZ_0$。置 $Q_S=Q\times_TS$。

因为 $X_i\simeq Z_i\times_TS$，遂得到\Term{isomorphism}{同构}
\[
d^{Z_*}\times_TS:X_1\xrightarrow\sim X_0\times_{Q_S}X_0.
\]
记 $\varphi$ 为其逆，并设 $\pi$ 是\Term{canonical}{典范}\Term{morphism}{态射}
\[
X_0\times_YX_0\longrightarrow X_0\times_{Q_S}X_0.
\]
\SourcePage{287}
则 $\varphi\circ\pi$ 是 $d_0\boxtimes d_1:X_1\xrightarrow\sim X_0\times_YX_0$ 的逆。因此，$X_*$ 定义的\Term{equivalence-relation}{等价关系}，即\Term{monomorphism}{单态射}
\[
\begin{tikzcd}
X_1\arrow[r,"d_0\boxtimes d_1","\sim"']&X_0\times_YX_0\arrow[r,hook]&X_0\times_SX_0
\end{tikzcd}
\]
与\Term{morphism}{态射} $q_S:X_0\to Q_S$ 所定义的\Term{equivalence-relation}{等价关系} $\mathscr R(q_S)$ 等同。后者\Term{faithfully-flat}{忠实平坦}且\Term{finite-presentation}{有限呈示}，因而为\Term{universal-effective-epimorphism}{普遍有效满态射}，所以 $\mathscr R(q_S)$ 的\Term{quotient}{商}是 $Q_S$（参见 IV，3.3.2）。故 $Y\simeq Q\times_TS$，从而 $Y\to S$ 和 $p$ \Term{finite-presentation}{有限呈示}。此外，由 IV，6.3.3，$(Y,p)$ 也是 $(d_0,d_1)$ 在 $(\mathrm{fppf})$ \Term{topology}{拓扑}的\Term{sheaf}{层}\Term{category}{范畴}中的\Term{coequalizer}{余核}。
\end{proof}

\begin{SGAresult}{命题}{9.1}
考虑\Term{scheme-morphism}{概形态射}
\[
X_0\xrightarrow{p}Y\xrightarrow{q}S,
\]
其中 $qp$ \Term{finite-type}{有限型}（分别\Term{finite-presentation}{有限呈示}），$p$ \Term{faithfully-flat}{忠实平坦}且\Term{finite-presentation}{有限呈示}。则 $q$ \Term{finite-type}{有限型}（分别\Term{finite-presentation}{有限呈示}）。\begingroup\renewcommand{\thefootnote}{*}\footnote{参见 EGA IV$_4$，17.7.5，其中给出了更一般的结果。}\addtocounter{footnote}{-1}\endgroup{}
\end{SGAresult}
\begin{proof}
因为 $p$ 为\Term{surjectivity}{满射}且 $qp$ \Term{quasi-compact}{拟紧}，$q$ \Term{quasi-compact}{拟紧}。因此可假设 $S,Y,X_0$ \Term{affine}{仿射}，对应的\Term{ring}{环}分别为 $A,B,C$。有 $B=\varinjlim B_i$，其中 $B_i$ 遍历 $B$ 的\Term{finite-type}{有限型} $A$-\Term{subalgebra}{子代数}。因为 $C$ 在 $B$ 上\Term{finite-presentation}{有限呈示}，存在指标 $i_0$、一个\Term{finite-presentation}{有限呈示} $B_{i_0}$-\Term{algebra}{代数} $C_{i_0}$ 及一个\Term{isomorphism}{同构} $C\simeq C_{i_0}\otimes_{B_{i_0}}B$；对 $i\geq i_0$ 置 $C_i=C_{i_0}\otimes_{B_{i_0}}B_i$，则 $C\simeq C_i\otimes_{B_i}B$。
\[
\begin{tikzcd}
B\arrow[r]&C\\
B_i\arrow[u]\arrow[r]&C_i\arrow[u]\\
A\arrow[u]&
\end{tikzcd}
\]
因为 $C$ 在 $B$ 上\Term{faithfully-flat}{忠实平坦}，由 EGA IV$_3$，11.2.6 和 8.10.5 (vi)，存在 $i_1\geq i_0$，使 $C_{i_1}$ 在 $B_{i_1}$ 上\Term{faithfully-flat}{忠实平坦}；所以对 $i\geq i_1$，$C_i$ 在 $B_i$ 上\Term{faithfully-flat}{忠实平坦}。此时，对 $i\geq i_1$，\Term{canonical}{典范}映射 $C_i\to C$ 是\Term{injectivity}{单射}，因为它由 $B_i\to B$ 经\Term{faithfully-flat}{忠实平坦}的底\Term{group-extension}{扩张}得到。

若 $C$ 在 $A$ 上\Term{finite-type}{有限型}，便存在指标 $j\geq i_1$ 使 $C_j=C$，故 $B_j=B$，因为 $C_j$ 在 $B_j$ 上\Term{faithfully-flat}{忠实平坦}。因此 $B$ 在 $A$ 上\Term{finite-type}{有限型}。

现在假设 $C$ 在 $A$ 上\Term{finite-presentation}{有限呈示}。由前述结果，$B$ 在 $A$ 上\Term{finite-type}{有限型}，故形如 $\overline B/I$，其中 $\overline B$ 是 $A$ 上\Term{finite}{有限}个不定元的\Term{polynomial-algebra}{多项式代数}，$I$ 是 $\overline B$ 的\Term{ideal}{理想}。$I$ 是其各\Term{finite-type}{有限型}\Term{subideal}{子理想} $I_\alpha$ 的并；故有 $B=\varinjlim B_\alpha$，其中 $B_\alpha=\overline B/I_\alpha$。与前面一样，存在指标 $\alpha_0$、一个\Term{finite-presentation}{有限呈示} $B_{\alpha_0}$-\Term{algebra}{代数} $C_{\alpha_0}$ 及一个\Term{isomorphism}{同构} $C\simeq C_{\alpha_0}\otimes_{B_{\alpha_0}}B$。对 $\alpha\geq\alpha_0$，仍置 $C_\alpha=C_{\alpha_0}\otimes_{B_{\alpha_0}}B_\alpha$，于是对 $\alpha\geq\alpha_0$，有 $C\simeq C_\alpha\otimes_{B_\alpha}B$。仍由 EGA IV$_3$，11.2.6
\SourcePage{288}
及 8.10.5 (vi)，与前面一样可知，当 $\alpha$ 充分大时，$C_\alpha$ 在 $B_\alpha$ 上\Term{faithfully-flat}{忠实平坦}。此时，映射 $C_\alpha\to C$（分别 $C_\alpha\to C_\beta$，$\beta\geq\alpha$）的\Term{kernel}{核}可与 $C_\alpha\otimes_{B_\alpha}(I/I_\alpha)$（分别 $C_\alpha\otimes_{B_\alpha}(I_\beta/I_\alpha)$）等同。

因为 $C_\alpha$ 和 $C$ 在 $A$ 上\Term{finite-presentation}{有限呈示}，且 $C_\alpha\to C$ 为\Term{surjectivity}{满射}，$C_\alpha\otimes_{B_\alpha}(I/I_\alpha)$ 是\Term{finite-type}{有限型}\Term{ideal}{理想}，\footnote{编者注：参见 EGA IV$_1$，1.4.4。}且是各\Term{ideal}{理想} $C_\alpha\otimes_{B_\alpha}(I_\beta/I_\alpha)$ 的并。所以当 $\beta$ 充分大时，$C_\alpha\otimes_{B_\alpha}(I_\beta/I_\alpha)=C_\alpha\otimes_{B_\alpha}(I/I_\alpha)$，从而 $I_\beta=I$（因为 $C_\alpha$ 在 $B_\alpha$ 上\Term{faithfully-flat}{忠实平坦}）；故 $B$ 在 $A$ 上\Term{finite-presentation}{有限呈示}。
\end{proof}

\SGAsection{10}{补充：对群概形取商}
以下第 10.2--10.4 节根据 M. Raynaud 的指点撰写，旨在将前面的\Term{theorem}{定理}应用于\Term{group-scheme}{群概形}\Term{action}{作用}的情形。为方便读者，先在第 10.1 节重录第 XVI 讲的陈述 2.1 至 2.3。

\textbf{10.1. 商的\Term{representability}{可表性}\Term{theorem}{定理}。}先“回顾”以下结果。
\begin{SGAresult}{\Term{theorem}{定理}}{10.1.1}
设 $S$ 是\Term{scheme}{概形}，$X,Y$ 是两个 $S$-\Term{scheme}{概形}，$f:X\to Y$ 是 $S$-\Term{morphism}{态射}。假设属于以下两种情形之一：

$\alpha$) \Term{morphism}{态射} $f$ \Term{locally-finite-presentation}{局部有限呈示}。

$\beta$) \Term{scheme}{概形} $S$ \Term{locally-noetherian}{局部诺特}，$X$ 在 $S$ 上\Term{locally-finite-type}{局部有限型}。

则下列条件\Term{equivalence}{等价}：

(i) 存在 $S$-\Term{scheme}{概形} $X'$ 及 $f$ 的\Term{factorization}{分解}
\[
f:X\xrightarrow{f'}X'\xrightarrow{f''}Y,
\]
其中 $f'$ 是\Term{faithfully-flat}{忠实平坦}且\Term{locally-finite-presentation}{局部有限呈示}的 $S$-\Term{morphism}{态射}，$f''$ 是\Term{monomorphism}{单态射}。

(ii) （第一）\Term{projection}{投影}
\[
p_1:X\times_YX\longrightarrow X
\]
是\Term{flat}{平坦}\Term{morphism}{态射}。

此外，若以上条件成立，$(X',f')$ 是 $X$ 对 $f$ 定义的\Term{equivalence-relation}{等价关系}的\Term{quotient}{商}（就 $(\mathrm{fppf})$ \Term{topology}{拓扑}而言），所以 (i) 中的\Term{factorization}{分解} $f=f''\circ f'$ 在\Term{isomorphism}{同构}意义下唯一。
\end{SGAresult}
$Y$ \Term{locally-noetherian}{局部诺特}、$X$ 在 $Y$ 上\Term{finite-type}{有限型}的情形已在 [Mur65]，\Term{theorem}{定理} 2 的推论 2 中处理。我们将说明如何化归于这一情形。

先作几点说明。
\SourcePage{289}
a) 蕴含 (i) $\Rightarrow$ (ii) 显然。事实上，第一\Term{projection}{投影}
\[
p'_1:X\times_{X'}X\longrightarrow X
\]
经由 $X\times_YX$ \Term{factorization}{分解}为
\[
p'_1:X\times_{X'}X\xrightarrow{u}X\times_YX\xrightarrow{p_1}X.
\]
因为 $f''$ 是\Term{monomorphism}{单态射}，$u$ 是\Term{isomorphism}{同构}；因为 $f'$ \Term{flat}{平坦}，$p'_1$ \Term{flat}{平坦}，故 $p_1$ \Term{flat}{平坦}。

b) 10.1.1 的断言在 $Y$ 上是\Term{local-scheme}{局部}的（因此在 $S$ 上也是\Term{local-scheme}{局部}的）；它们在 $X$ 上也是\Term{local-scheme}{局部}的，这容易从\Term{flat}{平坦}且\Term{locally-finite-presentation}{局部有限呈示}\Term{morphism}{态射}为\Term{open-morphism}{开态射}这一事实得出（EGA IV$_3$，11.3.1）。

c) 在 10.1.1 $\alpha$) 的假设下，由以上结果，我们化归到 $X,Y$ \Term{affine}{仿射}且 $f$ \Term{finite-presentation}{有限呈示}的情形。以 $Y$ 代替 $S$，可假设 $X,Y$ 在 $S$ 上\Term{finite-presentation}{有限呈示}。再利用 EGA IV$_3$，11.2.6 化归到 $S$ \Term{noetherian}{诺特}的情形。

d) 在 10.1.1 $\beta$) 的假设下，可假设 $S,X,Y$ \Term{affine}{仿射}，$S$ \Term{noetherian}{诺特}，$X$ 在 $S$ 上\Term{finite-type}{有限型}。将 $Y$ 看作在 $S$ 上\Term{finite-type}{有限型}的\Term{affine}{仿射}\Term{scheme}{概形} $Y_i$ 的\Term{filtered-inverse-limit}{滤过射影极限}。各\Term{scheme}{概形} $X\times_{Y_i}X$ 构成 $X\times_SX$ 的一个\Term{filtered}{滤过}递降的\Term{closed-subscheme}{闭子概形}族，其\Term{inverse-limit}{射影极限}为 $X\times_YX$。因 $X\times_SX$ \Term{noetherian}{诺特}，当 $i$ 充分大时，有 $X\times_{Y_i}X=X\times_YX$；因此，若 $f$ 满足 10.1.1 (ii) 的假设，则 $f_i:X\xrightarrow{f}Y\to Y_i$ 也满足。因为 $f$ 在 $X$ 上定义的\Term{equivalence-relation}{等价关系}与 $f_i$ 定义的相同，显然只须对 $f_i$ 证明 (ii) $\Rightarrow$ (i)，于是化归到 $Y$ 在 $S$ 上\Term{finite-type}{有限型}的情形。

\textit{在\Term{group-scheme}{群概形}上的应用。}设 $S$ 是\Term{scheme}{概形}，$G$ 是在 $S$ 上\Term{locally-finite-presentation}{局部有限呈示}的 $S$-\Term{group-scheme}{群概形}，它（从左）\Term{action}{作用}于 $S$-\Term{scheme}{概形} $X$。若 $X\to S$ 有\Term{section}{截面} $\xi$，回顾 $\operatorname{Stab}_G(\xi)$ 可由 $G$ 的一个\Term{subgroup-scheme}{子群概形}表示（参见 I，2.3.3）。

\begin{SGAresult}{\Term{theorem}{定理}}{10.1.2}
设 $S$ 是\Term{scheme}{概形}，$G$ 是在 $S$ 上\Term{locally-finite-presentation}{局部有限呈示}的 $S$-\Term{group-scheme}{群概形}，它\Term{action}{作用}于 $S$-\Term{scheme}{概形} $X$。

假设 $X\to S$ 有\Term{section}{截面} $\xi$，其在 $G$ 中的\Term{stabilizer}{稳定子} $H$ 在 $S$ 上\Term{flat}{平坦}。若满足下列假设之一：

a) $X$ 在 $S$ 上\Term{locally-finite-type}{局部有限型}，

b) $S$ \Term{locally-noetherian}{局部诺特}，

则 $(\mathrm{fppf})$ \Term{quotient-sheaf}{商层} $G/H$ 可由一个在 $S$ 上\Term{locally-finite-presentation}{局部有限呈示}的 $S$-\Term{scheme}{概形}表示，且 $S$-\Term{morphism}{态射}
\[
f:G\longrightarrow X,\qquad g\longmapsto g\cdot\xi
\]
\SourcePage{290}
\Term{factorization}{分解}为
\[
\begin{tikzcd}
G\arrow[d,"p"']\arrow[dr,"f"]&\\
G/H\arrow[r,hook,"i"']&X,
\end{tikzcd}
\]
其中 $p$ 是\Term{canonical}{典范}\Term{projection}{投影}，为\Term{faithfully-flat}{忠实平坦}\Term{locally-finite-presentation}{局部有限呈示}\Term{morphism}{态射}，$i$ 是\Term{monomorphism}{单态射}。
\end{SGAresult}
\begin{proof}
\Term{morphism}{态射} $f$ 使 $G$ 成为 $X$-\Term{scheme}{概形}。由 $\xi$ 的\Term{stabilizer}{稳定子}的定义，\Term{morphism}{态射}
\[
G\times_SH\longrightarrow G\times_XG,\qquad(g,h)\longmapsto(g,gh)
\]
是\Term{isomorphism}{同构}。因为 $H$ 在 $S$ 上\Term{flat}{平坦}，$G\times_SH$ 在 $G$ 上\Term{flat}{平坦}，所以第一\Term{projection}{投影} $p_1:G\times_SG\to G$ 是\Term{flat}{平坦}\Term{morphism}{态射}。另一方面，若 $X$ 在 $S$ 上\Term{locally-finite-type}{局部有限型}，则 $f$ \Term{locally-finite-presentation}{局部有限呈示}（EGA IV$_1$，1.4.3 (v)）；否则，假设 $S$ \Term{locally-noetherian}{局部诺特}。于是只须对\Term{morphism}{态射} $f$ 应用 10.1.1。还须证明 $G/H$ 在 $S$ 上\Term{locally-finite-presentation}{局部有限呈示}，但这立即由 9.1 得出。
\end{proof}

\begin{SGAresult}{推论}{10.1.3}
设 $S$ 是\Term{scheme}{概形}，$u:G\to H$ 是 $S$-\Term{group-scheme}{群概形}的\Term{morphism}{态射}。假设 $G$ 在 $S$ 上\Term{locally-finite-presentation}{局部有限呈示}，且 $H$ 在 $S$ 上\Term{locally-finite-type}{局部有限型}，或者 $S$ \Term{locally-noetherian}{局部诺特}。

那么，若 $K=\operatorname{Ker}(u)$ 在 $S$ 上\Term{flat}{平坦}，\Term{quotient-group}{商群} $G/K$ 可由一个在 $S$ 上\Term{locally-finite-presentation}{局部有限呈示}的 $S$-\Term{group-scheme}{群概形}表示，而且 $u$ \Term{factorization}{分解}为
\[
\begin{tikzcd}
G\arrow[r,"u"]\arrow[dr,"p"']&H\\
&G/K\arrow[u,hook,"i"']
\end{tikzcd}
\]
其中 $p$ 为\Term{canonical}{典范}\Term{projection}{投影}，$i$ 为\Term{monomorphism}{单态射}。
\end{SGAresult}
\begin{proof}
应用 10.1.2，取 $X=H$，$\xi$ 为 $H$ 的\Term{unit-section}{单位截面}。
\end{proof}

\textbf{10.2. \Term{diagonal}{对角线}的\Term{stabilizer}{稳定子}。}设 $S$ 是\Term{noetherian}{诺特}\Term{scheme}{概形}，$X$ 是\Term{finite-type}{有限型} $S$-\Term{scheme}{概形}，$G$ 是\Term{flat}{平坦}\Term{finite-type}{有限型} $S$-\Term{group-scheme}{群概形}，从左\Term{action}{作用}于 $X$，即给定一个 $S$-\Term{action}{作用} $d_0:G\times_SX\to X$。记 $d_1:G\times_SX\to X$ 为向第二因子的\Term{projection}{投影}。按第 2 节 a)，有\Term{groupoid}{群胚}
\[
\begin{tikzcd}[column sep=large]
G\times_SG\times_SX\arrow[r,shift left=1.2ex,"{\operatorname{pr}_{2,3}}"]\arrow[r,"\mu\times X" description]\arrow[r,shift right=1.2ex,"G\times d_0"']&G\times_SX\arrow[r,shift left=.6ex,"d_1"]\arrow[r,shift right=.6ex,"d_0"']&X,
\end{tikzcd}
\]
回顾其\Term{coequalizer}{余核}若存在，记作 $G\backslash X$。
\SourcePage{291}
\begin{SGAresult}{定义}{10.2.1}
以 $F\subset G\times_SX$ 表示\Term{diagonal-section}{对角截面}的\Term{stabilizer}{稳定子}，即由\Term{fiber-product}{纤维积}
\[
\begin{tikzcd}
F\arrow[r]\arrow[d]&X\arrow[d,"\Delta"]\\
G\times_SX\arrow[r,"{(d_0,d_1)}"']&X\times_SX
\end{tikzcd}
\]
定义的 $X$-\Term{scheme}{概形}。
\end{SGAresult}
则 $F$ 是 $G\times_SX$ 的一个子 $X$-\Term{group-scheme}{群概形}。因 $G\times_SX$ 在\Term{noetherian}{诺特}\Term{scheme}{概形} $S$ 上\Term{finite-type}{有限型}，因而\Term{noetherian}{诺特}，所以 $F$ 在 $S$ 上和 $X$ 上都\Term{finite-type}{有限型}（EGA I，6.3.5 和 6.3.6）。此外，若 $X\to S$ \Term{separated}{分离}，则 $F$ 是 $G\times_SX$ 的一个闭子 $X$-\Term{group-scheme}{群概形}。

回顾，若\Term{morphism}{态射}
\[
G\times_SX\xrightarrow{(d_0,d_1)}X\times_SX
\]
是\Term{monomorphism}{单态射}，则称 $G$ 在 $X$ 上\Term{free-action}{自由作用}（参见第 III 讲 3.2.1）。这\Term{equivalence}{等价}于 $F$ 是以 $X$ 为底的\Term{trivial-group-scheme}{平凡群概形}。

\textbf{10.3. $F$ 在 $X$ 上\Term{quasi-finite}{拟有限}的情形。}因 $F$ 在 $X$ 上\Term{finite-type}{有限型}，它在 $X$ 上\Term{quasi-finite}{拟有限}，当且仅当 $X$ 的各\Term{geometric-point}{几何点}的\Term{stabilizer}{固定子}\Term{finite}{有限}。
\begin{SGAresult}{\Term{theorem}{定理}}{10.3.1}
\footnote{编者注：这里也存在 $X$ 的一个满足\Term{theorem}{定理}结论的最大\Term{open-subset}{开子集} $U$，参见编者注 (44)。}
在 10.2 的假设下，假设 $F$ 在 $X$ 上\Term{quasi-finite}{拟有限}。则 $X$ 中存在一个\Term{dense}{稠密}且 $G$-\Term{saturated}{饱和}的\Term{open-subset}{开子集} $U$，满足下列性质：

(i) 在 $(\mathrm{Sch}/S)$ 中，\Term{coequalizer}{余核} $V=G\backslash U$ 存在；此外，\Term{scheme}{概形} $V$ 是\Term{ringed-space}{赋环空间}\Term{category}{范畴}中的\Term{quotient}{商}。

(ii) $p:U\to V$ 为\Term{surjectivity}{满射}、\Term{open-morphism}{开态射}且\Term{finite-presentation}{有限呈示}。

(iii) $V$ 在 $S$ 上\Term{finite-presentation}{有限呈示}。

(iv) \Term{morphism}{态射} $G\times_SU\to U\times_VU$，$(g,x)\mapsto(gx,x)$，为\Term{surjectivity}{满射}。

(v) 再假设 $G$ 在 $X$ 上\Term{free-action}{自由作用}。则 $U\to V$ 是一个对 $(\mathrm{fppf})$ \Term{topology}{拓扑}\Term{locally-trivial}{局部平凡}的（左）$G$-\Term{torsor}{挠子}。特别地，$U\to V$ \Term{faithfully-flat}{忠实平坦}。\footnote{编者注：若再假设 $G$ 为\Term{reductive-group-scheme}{约化} $S$-\Term{group-scheme}{群概形}，且 $G$ 在 $X$ 上的（自由）\Term{action}{作用}可\Term{linearizable-action}{线性化}，则已知 $G\backslash X$ \Term{representable}{可表}，且 $X\to G\backslash X$ 是一个（左）$G$-\Term{torsor}{挠子}。这是 Raynaud 和 Seshadri 的结果的推论，可见文章 [CTS79]（命题 6.11）。}
\end{SGAresult}
\begin{proof}
已假设\Term{morphism}{态射} $G\times_SX\to X\times_SX$，$(g,x)\mapsto(gx,x)$，是\Term{quasi-finite}{拟有限}的。因此可对 $(X,G)$ 所定义的\Term{groupoid}{群胚}应用\Term{theorem}{定理} 8.1。于是存在\Term{dense}{稠密}\Term{saturated}{饱和}\Term{open-subset}{开子集} $U\subset X$，使\Term{quotient}{商} $G\backslash U$ 存在；它满足性质 (i)、(ii)、(iii)。
\SourcePage{292}
为证明 (iv)，回顾 $G$ 在 $X$ 上\Term{free-action}{自由作用}，当且仅当 $(d_0,d_1)$ 是\Term{equivalence-pair}{等价对}（III，3.2.1）。此时，\Term{theorem}{定理} 8.1 (iv) 表明，\Term{morphism}{态射} $G\times_SU\to U\times_VU$ 是\Term{isomorphism}{同构}，且 $p$ \Term{faithfully-flat}{忠实平坦}且\Term{finite-presentation}{有限呈示}。因此，$U$ 是以 $V$ 为底、对 $(\mathrm{fppf})$ \Term{topology}{拓扑}\Term{locally-trivial}{局部平凡}的 $G$-\Term{torsor}{挠子}。
\end{proof}

\textbf{10.4. $F$ 在 $X$ 上\Term{flat}{平坦}的情形。}记
\[
d=(d_0,d_1):G\times_SX\longrightarrow X\times_SX
\]
为\Term{morphism}{态射} $d(g,x)=(gx,x)$。回顾，与 $(X,G)$ 相伴的\Term{equivalence-relation}{等价关系}的\Term{sheaf-theoretic-graph}{层论图像} $\widetilde\Gamma$，是 $X\times_SX$ 的 $S$-子 $(\mathrm{fppf})$ \Term{sheaf}{层}，即 $(d_0,d_1)$ 的像。它是与\Term{graph-functor}{图像函子}
\[
T\longmapsto\Gamma(T)=\{(x_0,x_1)\in X(T)\times X(T)\mid x_0\in G(T)x_1\}
\]
相伴的 $(\mathrm{fppf})$ \Term{sheaf}{层}。

置 $G_X=G\times_SX$。对任意 $S$-\Term{scheme}{概形} $T$，有\Term{surjectivity}{满射}
\[
G_X(T)\longrightarrow\Gamma(T),\qquad(g,x)\longmapsto(gx,x),
\]
它诱导\Term{bijection}{双射}
\[
\varphi(T):G_X(T)/F(T)\xrightarrow\sim\Gamma(T);
\]
事实上，若 $(g,x),(g',x')\in G_X(T)$ 满足 $(gx,x)=(g'x',x')$，则 $x'=x$ 且 $g^{-1}g'x=x$，故 $(g^{-1}g'x,x)\in F(T)$，且 $(g,x)$ 与 $(g',x)$ 在 $G_X(T)/F_X(T)$ 中有相同的像。

由定义（参见 IV，4.4.1 (ii) 或 5.2.1 的证明），\Term{quotient-sheaf}{商层} $G_X/F$ 是与\Term{presheaf}{预层}
\[
T\longmapsto G_X(T)/F(T)\simeq\Gamma(T)
\]
相伴的 $(\mathrm{fppf})$ \Term{sheaf}{层}。因此有\Term{sheaf-isomorphism}{层同构} $\varphi:G_X/F\to\widetilde\Gamma$。

\begin{SGAresult}{\Term{theorem}{定理}}{10.4.1}
\footnote{编者注：这是 [Ray67b] 的\Term{theorem}{定理} 3 的 (2)。该短札还概述了\Term{theorem}{定理} 10.1.1 的另一个证明。}
在 10.2 的假设下：

a) $\widetilde\Gamma$ \Term{representable}{可表}，当且仅当 $F$ 在 $X$ 上\Term{flat}{平坦}。

b) 假设 $F$ 在 $X$ 上\Term{flat}{平坦}。则 $d_1,d_0$ 所诱导的\Term{morphism}{态射}
\[
\begin{tikzcd}
G_X/F\arrow[r,shift left=.6ex,"\overline d_1"]\arrow[r,shift right=.6ex,"\overline d_0"']&X
\end{tikzcd}
\]
\Term{faithfully-flat}{忠实平坦}且\Term{finite-presentation}{有限呈示}。
\end{SGAresult}
\SourcePage{293}
\begin{proof}
a) 的证明。假设 $(\mathrm{fppf})$ \Term{sheaf}{层} $G_X/F$ 可由一个 $X$-\Term{scheme}{概形} $Y$ 表示。那么，由 IV，6.3.3，$p:G_X\to Y$ \Term{faithfully-flat}{忠实平坦}且\Term{locally-finite-presentation}{局部有限呈示}，且下图的第二个方块是\Term{cartesian}{笛卡尔}的：
\[
\begin{tikzcd}[column sep=large]
F\arrow[r]\arrow[d]&F\times_XG_X\arrow[r]\arrow[d]&G_X\arrow[d,"p"]\\
X\arrow[r,"e_X"']&G_X\arrow[r]&Y,
\end{tikzcd}
\]
第一个方块由\Term{unit-section}{单位截面} $e_X:X\to G_X$ 的\Term{base-change}{基变换}得到。因为 $p$ \Term{faithfully-flat}{忠实平坦}且\Term{locally-finite-presentation}{局部有限呈示}，$F\to X$ 也如此。

反过来，假设 $F$ 在 $X$ 上\Term{flat}{平坦}。置 $X^2=X\times_SX$。\Term{morphism}{态射} $d:G_X\to X^2$ 可用来构造\Term{fiber-product}{纤维积}
\[
\begin{tikzcd}
G_X\times_{X^2}G_X\arrow[r]\arrow[d]&G_X\arrow[d]\\
G_X\arrow[r]&X^2.
\end{tikzcd}
\]
则\Term{morphism}{态射} $G_X\times_{X^2}G_X\to X^2$ 是 $X^2$ 上的一个 $F\times_XX^2$-\Term{torsor}{挠子}，故\Term{flat}{平坦}且\Term{finite-type}{有限型}（因为 $F$ 如此）。由\Term{theorem}{定理} 10.1.1，\Term{morphism}{态射} $d$ 唯一\Term{factorization}{分解}为
\[
G_X\xrightarrow{\psi}Y\xrightarrow{\tau}X\times_SX,
\]
其中 $\psi$ \Term{faithfully-flat}{忠实平坦}（且\Term{finite-type}{有限型}），$\tau$ 是\Term{scheme}{概形}的\Term{monomorphism}{单态射}。

于是\Term{sheaf-morphism}{层态射} $\psi:G_X\to Y$ 是 $F$-\Term{invariant}{不变}的，因而得到\Term{sheaf-morphism}{层态射} $\overline\psi:G_X/F\to Y$。另一方面，由于 $\psi$ \Term{faithfully-flat}{忠实平坦}（且\Term{finite-type}{有限型}），\Term{sheaf}{层}的\Term{monomorphism}{单态射} $\tau$ 经由 $d$ 的\Term{image-sheaf}{像层}，即 $\widetilde\Gamma$ \Term{factorization}{分解}。因此，\Term{sheaf-isomorphism}{层同构} $G_X/F\simeq\widetilde\Gamma$ 经由\Term{monomorphism}{单态射} $Y\to\widetilde\Gamma$ \Term{factorization}{分解}。所以 $Y$ 表示 $G_X/F$。

b) 的证明。假设 $F$ 在 $X$ 上\Term{flat}{平坦}。由 a) 及其证明，$G_X/F$ \Term{representable}{可表}，\Term{morphism}{态射} $p:G_X\to G_X/F$ \Term{faithfully-flat}{忠实平坦}且\Term{finite-presentation}{有限呈示}。另一方面，由假设，\Term{morphism}{态射} $d_i:G_X\to X$（$i=0,1$）\Term{faithfully-flat}{忠实平坦}且\Term{finite-presentation}{有限呈示}。因 $d_i=\overline d_i\circ p$，由 EGA IV$_2$，2.2.13 (iii) 及 EGA IV$_3$，11.3.16，$\overline d_i$ \Term{faithfully-flat}{忠实平坦}且\Term{finite-presentation}{有限呈示}。
\end{proof}

\begin{SGAresult}{\Term{theorem}{定理}}{10.4.2}
\footnote{编者注：这里也存在 $X$ 中满足\Term{theorem}{定理}结论的最大\Term{open-subset}{开子集} $U$；此外，$X$ 中\Term{codimension}{余维}为 $1$ 的点 $x$ 属于 $U$，当且仅当\Term{morphism}{态射} $(G_X/F)\times_X\operatorname{Spec}(\mathcal O_{X,x})\to\operatorname{Spec}(\mathcal O_{X,x})\times_S\operatorname{Spec}(\mathcal O_{X,x})$ 是\Term{closed-immersion}{闭浸入}，参见编者注 (44)。}
在 10.2 的假设下，假设 $F$ 在 $X$ 上\Term{flat}{平坦}。则 $X$ 中存在\Term{dense}{稠密}\Term{saturated}{饱和}\Term{open-subset}{开子集} $U$，使 $(\mathrm{fppf})$ \Term{quotient}{商} $V=G\backslash U$ 是\Term{finite-type}{有限型} $S$-\Term{scheme}{概形}，且 $U\to V$ \Term{faithfully-flat}{忠实平坦}且\Term{finite-presentation}{有限呈示}。
\end{SGAresult}
\SourcePage{294}
\begin{proof}
\Term{theorem}{定理} 10.4.1 表明，$G_X/F\simeq\widetilde\Gamma$ \Term{representable}{可表}。于是 $(\mathrm{fppf})$ \Term{sheaf}{层} $G\backslash X$ 与
\[
\begin{tikzcd}
G_X/F\arrow[r,shift left=.6ex,"\overline d_1"]\arrow[r,shift right=.6ex,"\overline d_0"']&X
\end{tikzcd}
\]
的\Term{quotient-sheaf}{商层}等同。由前述结果，$\overline d_i:G_X/F\to X$ \Term{faithfully-flat}{忠实平坦}且\Term{finite-presentation}{有限呈示}（$i=0,1$），而\Term{morphism}{态射}
\[
\begin{tikzcd}
G_X/F\arrow[r,"\sim"]&\widetilde\Gamma\arrow[r,hook]&X\times_SX
\end{tikzcd}
\]
是\Term{monomorphism}{单态射}，即 $(\overline d_0,\overline d_1)$ 是\Term{equivalence-pair}{等价对}。因此可应用\Term{theorem}{定理} 8.1。故存在 $X$ 的\Term{dense}{稠密}\Term{saturated}{饱和}\Term{open-subset}{开子集} $U$，使 $(\mathrm{fppf})$ \Term{quotient}{商} $V=G\backslash U$ 是\Term{finite-type}{有限型} $S$-\Term{scheme}{概形}，且 $U\to V$ \Term{faithfully-flat}{忠实平坦}且\Term{finite-presentation}{有限呈示}。
\end{proof}
结合\Term{generic-flatness-theorem}{泛平坦性定理}（EGA IV$_2$，6.9.3），得到：
\begin{SGAresult}{推论}{10.4.3}
在 10.2 的假设下，假设 $X$ \Term{reduced}{约化}。则 $X$ 中存在\Term{dense}{稠密}\Term{saturated}{饱和}\Term{open-subset}{开子集} $U$，使 $(\mathrm{fppf})$ \Term{quotient}{商} $G\backslash U$ 是\Term{finite-type}{有限型} $S$-\Term{scheme}{概形}，且 $U\to G\backslash U$ \Term{faithfully-flat}{忠实平坦}且\Term{finite-presentation}{有限呈示}。
\end{SGAresult}

\textbf{参考文献}\footnote{编者注：本讲引用的补充文献。}
\begin{description}
\item[{[AK80]}] A. B. Altman、S. L. Kleiman，\textit{Compactifying the Picard Scheme}（\Term{picard-scheme}{Picard 概形}的\Term{compactification}{紧化}），Adv. Math. \textbf{35} (1980)，50--112。
\item[{[An73]}] S. Anantharaman，\textit{Schémas en groupes, espaces homogènes et espaces algébriques sur une base de dimension 1}（一维底上的\Term{group-scheme}{群概形}、\Term{homogeneous-space}{齐性空间}与\Term{algebraic-space}{代数空间}），Mém. Soc. Math. France \textbf{33} (1973)，5--79。
\item[{[BLR90]}] S. Bosch、W. Lütkebohmert、M. Raynaud，\textit{Néron models}（\Term{neron-model}{Néron 模型}），Springer-Verlag，1990。
\item[{[CTS79]}] J.-L. Colliot-Thélène、J.-J. Sansuc，\textit{Fibrés quadratiques et composantes connexes réelles}（二次丛与实\Term{connected-component}{连通分支}），Math. Ann. \textbf{244} (1979)，105--134。
\item[{[DG70]}] M. Demazure、P. Gabriel，\textit{Groupes algébriques}（\Term{algebraic-group}{代数群}），Masson \& North-Holland，1970。
\item[{[DR81]}] J. Dixmier、M. Raynaud，\textit{Sur le quotient d'une variété algébrique par un groupe algébrique}（关于\Term{algebraic-variety}{代数簇}对\Term{algebraic-group}{代数群}的商），第 327--344 页，载于 \textit{Mathematical Analysis and Applications}（数学分析及其应用，L. Schwartz 六十五岁生日纪念，L. Nachbin 编），Adv. Math. Suppl. Stud.，第 7A 卷，1981。
\item[{[Fe03]}] D. Ferrand，\textit{Conducteur, descente et pincement}（\Term{derivation}{导子}、\Term{descent}{下降}与夹合），Bull. Soc. Math. France \textbf{131} (2003)，第 4 期，553--585。
\item[{[Hi62]}] H. Hironaka，\textit{An example of a non-Kählerian complex analytic deformation of Kählerian complex structures}（Kähler 复结构的一个非 Kähler 复解析\Term{deformation}{形变}例），Ann. of Math. \textbf{75} (1962)，第 1 期，190--208。
\end{description}
\SourcePage{295}
\begin{description}
\item[{[KM97]}] S. Keel、S. Mori，\textit{Quotient by groupoids}（对\Term{groupoid}{群胚}取商），Ann. of Math. \textbf{145} (1997)，第 1 期，193--213。
\item[{[Ko97]}] J. Kollár，\textit{Quotient spaces modulo algebraic groups}（模\Term{algebraic-group}{代数群}的\Term{quotient-space}{商空间}），Ann. of Math. \textbf{145} (1997)，第 1 期，33--79。
\item[{[Ko08]}] J. Kollár，\textit{Quotients by finite equivalence relations}（对\Term{finite}{有限}\Term{equivalence-relation}{等价关系}取商），arXiv: 0812.3608。
\item[{[Mum65]}] D. Mumford，\textit{Geometric invariant theory}（几何\Term{invariants}{不变量}理论），Springer-Verlag，1965；第 2 版（分别第 3 版）与 J. Fogarty（分别再与 F. Kirwan）合著，1982（分别 1994）。
\item[{[Mur65]}] J. P. Murre，\textit{Representation of unramified functors. Applications (according to unpublished results of A. Grothendieck)}（\Term{unramified}{非分歧}\Term{functor}{函子}的表示及应用，依据 A. Grothendieck 未发表的结果），Sém. Bourbaki，第 9 卷，第 294 讲 (1965)，Soc. Math. France，1995。
\item[{[Ray67a]}] M. Raynaud，\textit{Passage au quotient par une relation d'équivalence plate}（对\Term{flat}{平坦}\Term{equivalence-relation}{等价关系}取商），第 78--85 页，载于 \textit{Proc. Conf. Local Fields (Driebergen)}（\Term{local-scheme}{局部}域会议论文集，Driebergen，T. A. Springer 编），Springer-Verlag，1967。
\item[{[Ray67b]}] M. Raynaud，\textit{Sur le passage au quotient par un groupoïde plat}（关于对\Term{flat}{平坦}\Term{groupoid}{群胚}取商），C. R. Acad. Sci. Paris (Sér. A) \textbf{265} (1967)，384--387。
\end{description}
\SourcePage{296}
% 原 PDF 第 296 页为空白页。
